\subsection*{What automated reviewers catch in planted-flaw analyses}

We first asked how many planted flaws a reviewer configuration identifies, and whether a panel of three distinct lenses identifies more than a single lens sampled three times. All 1,296 reviewer calls (18 items $\times$ 8 calls $\times$ 3 repeats $\times$ 3 model configurations) returned a parseable critique block (one call needed a retry), with a mean of 17.3 critiques per call. The two judges agreed closely: Cohen's $\kappa$ was 0.93 (95\% CI 0.90--0.95) for the critique-level labels of 22,398 critiques and 0.84 (0.78--0.90) for whether each planted flaw was detected by each condition. Table~\ref{table_benchmark} summarizes the main outcomes.

\textbf{Several calls catch more than one; distinct lenses do not add recall at a matched number of calls.} A single rigour reviewer identified 81\% of planted flaws (pooled over configurations; 95\% CI 72--88\%), and a single combined-checklist reviewer 86\% (76--94\%) (Fig~\ref{fig2}A). Every three-call condition identified more: 90\% (83--96\%) for three rigour samples, 91\% (84--97\%) for three combined-checklist samples and 91\% (86--97\%) for the three-lens panel. Three samples of the same reviewer exceeded one sample by 9.5 percentage points for the rigour lens (95\% CI 5.3--13.7) and by 5.3 points for the combined-checklist reviewer (2.1--9.3; both Holm-adjusted $p=0.001$); because each three-call condition contains the single call it is compared with, these contrasts measure what the additional calls add (Materials and methods). The panel exceeded a single rigour reviewer by 10.9 points (6.3--15.7; $p_\mathrm{Holm}=0.001$), but did not exceed the self-ensembles with the same number of calls: the differences from three rigour samples and three combined samples were 1.4 (--1.9 to 4.4; $p_\mathrm{Holm}=0.82$) and 0.2 points (--1.6 to 2.5; $p_\mathrm{Holm}=0.98$). The result was robust to the judge: under the second judge, the panel exceeded a single rigour reviewer by 8.8 points and differed from the two self-ensembles by +0.5 and --0.5 points. In no configuration did the panel significantly exceed either self-ensemble. Its advantage over three rigour samples was largest for the low-effort configuration (+6.3 points, $p_\mathrm{Holm}=0.11$), whereas for gpt-5.5 three rigour samples exceeded the panel under both judges (2.8 points, $p_\mathrm{Holm}=0.14$, and 3.5 points, $p_\mathrm{Holm}=0.03$). For some models, therefore, spending two of the three calls on other lenses rather than on further rigour samples can cost a little recall. Overall, the gain of the panel over one reviewer is a gain from making three calls, not from making them through different lenses. This agrees with reports that budget-matched sampling and aggregation match multi-agent debate on reasoning tasks~\cite{wang2023selfconsistency,li2024moreagents,smit2024mad} and suggests the same holds for critique.

\textbf{The lenses largely overlap.} Of the 432 planted flaw instances reviewed by the panel, 359 (83\%) were identified by at least two of the three lenses, 10 (2.3\%) by the rigour lens alone, 22 (5.1\%) by the adversarial lens alone, 4 (0.9\%) by the biological-plausibility lens alone and 37 (8.6\%) by none; of the flaws the panel detected, 91\% were raised by two or more lenses. Part of this overlap follows from the prompts, which share some items by design (Materials and methods); the remainder suggests, but does not establish, that the models we tested apply much of a general rigour checklist even when prompted for adversarial or biological critique. Our design cannot separate these two sources, but either way lens specialisation changed the emphasis of the critiques more than their coverage.

\textbf{Technical errors about the model itself are the blind spot.} The panel's recall was high for statistical (98\%), confounding (97\%), causal-logic (94\%) and biological flaws (100\%), and every three-call condition reached 93--100\% in each of these families; single reviewers were less dependable there (82--97\% per family), with the lowest recall for circular validation (64--67\%) and, for the single rigour reviewer, effect-size neglect (61\%). Recall was low for technical flaws in every condition (61\% for the panel; Fig~\ref{fig2}B). The deficit is specific: false statements about attention mechanics were caught almost always---the three planted instances of ``dot-product attention is symmetric, so $A_{ij}=A_{ji}$'' were identified by 96\% of single rigour reviewers and by every multi-call condition---whereas false statements about the specific model or its tokenization (for example, that Geneformer uses binned expression values, or an incorrect layer count for a named model version) had a recall of only 17--28\% across conditions (four planted instances; 95\% CIs from 0\% to 33--56\%), with no consistent gain from additional calls or lenses. All detections were of the two statements about Geneformer; the two statements about scGPT (that it orders genes by expression rank, and a wrong layer, head and dimension count for its whole-human checkpoint) were never identified in any condition. Together with the near-universal detection of the attention-symmetry error, which no lens checklist names, this pattern suggests that reviewers check the logic of an analysis against general knowledge but do not reliably check its premises about the particular model being analysed. This is an inference from four flaws, and two contributing factors cannot be separated from it: the reviewers ran without tools and so could not look up model facts, and the rigour checklist mentions Geneformer's rank-value encoding but no fact about scGPT. Whatever the cause, such premises are better verified by checks that do not rely on the reviewer's recall, for example by checking architecture and tokenization claims against the model configuration in executed code.

\textbf{The panel is less noisy.} Reviewers raised far more concerns than there were planted flaws: across conditions only 19--29\% of critiques of flawed items matched a planted flaw, 56--64\% were specific but unplanted concerns, 11--22\% were generic and 1--2\% were judged incorrect. The panel had the highest share of critiques matching planted flaws (29\% vs 19--21\% for the other conditions) and the lowest share of generic critiques (11\% vs 15--22\%). On clean items, which contain no planted flaw, the panel raised a mean of 11.4 high or critical critiques per item, against 19.4 for both self-ensembles and 6.0--6.4 for single reviewers (Fig~\ref{fig2}C); the difference held with TF-IDF merging at $\tau=0.075$ (11.8 vs 18.2 and 21.8; S2 Table). The panel also produced fewer output tokens than either self-ensemble (6,100 vs 7,400 and 10,900 per item; Fig~\ref{fig2}D) and took less summed call time (141 vs 170 and 227\,s per item). Input tokens, dominated by the three copies of the write-up under review, were similar (33,100 vs 32,700 and 37,900 per item), so total tokens were only 2\% lower than for three rigour samples and 20\% lower than for three combined-checklist samples. At equal recall, distinct lenses thus buy a more focused critique list with fewer output tokens and shorter calls; against a single-lens self-ensemble the total-token saving is small.

\textbf{The letter grade does not certify a clean analysis.} Severity scores separated flawed from clean items (AUROC 0.84 for the panel, 95\% CI 0.79--0.90, and 0.74--0.85 for the other conditions; the flawed version scored higher than its clean counterpart in 98\% of matched comparisons). The panel's grade had the highest AUROC (0.89, 0.82--0.95, vs 0.74--0.82, with overlapping intervals), but this advantage depended on the pre-specified lexical merge, which left the panel's critiques almost entirely unmerged while merging and escalating some near-identical critiques within the self-ensembles: with TF-IDF merging at $\tau=0.075$, the grade of every multi-call condition had an AUROC of 0.50 (next section). Yet clean items almost never received a passing grade: 268 of 270 clean-item grades across conditions were D or F, and the remaining two were C. Because every reviewer finds some high-severity concern in any analysis, the grade is at most informative as a ranking, not as an absolute criterion of acceptability, which has direct consequences for stopping rules (below).
